\chapter{Shoulder Pain}

\section*{Definition}
Pain or discomfort localized to the shoulder joint and surrounding structures, arising from the glenohumeral joint, acromioclavicular joint, rotator cuff, or referred from cervical spine or viscera.

\section*{Pathophysiology}
The shoulder's complex anatomy allows extensive range of motion but predisposes to instability and impingement. Rotator cuff muscles (supraspinatus, infraspinatus, teres minor, subscapularis) stabilize the glenohumeral joint; their dysfunction leads to impingement, tendinopathy, or tears. Referred pain from the diaphragm, gallbladder, or heart may present as shoulder pain via phrenic nerve connections.

\section*{Causes}
\begin{itemize}
  \item Rotator cuff tendinopathy or tear
  \item Frozen shoulder (adhesive capsulitis)
  \item Shoulder impingement syndrome
  \item Acromioclavicular joint arthritis or separation
  \item Glenohumeral osteoarthritis
  \item Bicipital tendinopathy
  \item Labral tear (SLAP tear)
  \item Cervical radiculopathy (referred from neck)
  \item Fracture (proximal humerus, clavicle)
  \item Referred cardiac pain (left shoulder — angina or MI)
\end{itemize}

\section*{When to See a Doctor}
See your doctor if:
\begin{itemize}
  \item Severe or worsening symptoms
  \item Acute shoulder pain after injury, inability to raise arm, fever, or chest pain radiating to shoulder
  \item Accompanied by other concerning signs
\end{itemize}

\section*{Related Symptoms}
\begin{itemize}
  \item Neck pain
  \item Arm weakness
  \item Limited range of motion
  \item Night pain
  \item Arm numbness
\end{itemize}
\textbf{Important:} If this symptom persists, your doctor will check for serious underlying causes.

\section*{Self-Care / Home Management}
\begin{itemize}
  \item Rest and avoid activities that make it worse.
  \item Maintain adequate hydration and a balanced diet.
  \item Over-the-counter remedies may provide symptomatic relief (consult a pharmacist or doctor).
  \item Monitor symptoms and seek professional advice if they persist or worsen.
  \item Avoid self-medicating with prescription medications without medical guidance.
\end{itemize}

\section*{How Your Doctor Will Evaluate You}
\begin{itemize}
  \item A thorough history and physical examination are the first steps in evaluation.
  \item Laboratory tests (blood work, urinalysis) may be ordered based on clinical suspicion.
  \item Imaging studies (X-ray, ultrasound, CT, MRI) may be indicated when structural pathology is suspected.
  \item Specialist referral is arranged when the diagnosis remains uncertain or requires specific expertise.
\end{itemize}

\section*{Treatment Options}
\begin{itemize}
  \item Treatment targets the underlying cause identified through diagnostic evaluation.
  \item Supportive care and symptom management are provided while awaiting definitive therapy.
  \item Medications, lifestyle modifications, and physical therapies may be prescribed as appropriate.
  \item Follow-up ensures treatment effectiveness and allows timely adjustments to the care plan.
\end{itemize}

\clearpage